\documentclass[11pt]{article}
\usepackage{apchem}
\doctitle{Unit 4 Content Guide}
\docsubtitle{A Comprehensive Review of Chemical Reactions, Topics 4.1--4.9}
\docstatus{Study draft}
\begin{document}
\titleblock
\begin{abstract}
This guide covers all of Unit 4, Chemical Reactions, in ten sections: one for each of Topics 4.1--4.9 in course order, plus a final section on lab skills and free-response technique. Each section gives the definitions, equations and particle-level reasoning the topic needs, at least one worked example, the phrasing that earns free-response points and the mistakes that cost them. The worked examples use their own reactions and numbers, so none of them give away an item in \emph{Unit 4 Practice Set A} or \emph{Unit 4 Practice Set B}; the answer keys of both sets point back to the section numbers here. Plan on about 70 minutes for a full read with a calculator out.
\end{abstract}
\leadline{Usage}{first pass, read Sections 1--10 in order and redo every worked example. Then take \emph{Unit 4 Practice Set A} cold; reread the section behind every miss, and take \emph{Unit 4 Practice Set B} two or three days later. The last page is a one-page summary for the night before the unit test. Unit 4 is 7--9\% of the AP multiple-choice section, but stoichiometry, net ionic equations and titration return in nearly every later unit.}

\section{Introduction for reactions (Topic 4.1)}\label{t:4.1}
A chemical reaction rearranges atoms into new substances. A balanced equation shows that atoms (and therefore mass) are conserved: the same number of each kind of atom appears on both sides.
\begin{itemize}
\item Evidence of a reaction: a precipitate forms, a gas is released (not by boiling), a color changes, or energy is released or absorbed by a change in substances. None of these is proof on its own; a new substance is the test.
\item Balancing: change coefficients, never subscripts. Balance elements that appear in only one compound on each side first; leave free elements ($\ce{O2}$, $\ce{H2}$, metals) and H and O for last.
\item Coefficients give the ratio of particles, and so the ratio of moles. They do not give a ratio of masses.
\end{itemize}
\worked{1.1} Balance $\ce{C2H6 + O2 -> CO2 + H2O}$. C: 2 $\ce{CO2}$; H: 3 $\ce{H2O}$; O on the right $4 + 3 = 7$, so $\tfrac72\ce{O2}$. Double everything: $\ce{2C2H6 + 7O2 -> 4CO2 + 6H2O}$.

\section{Net ionic equations (Topic 4.2)}\label{t:4.2}
\begin{enumerate}
\item Write the balanced molecular equation with states.
\item Split every \emph{strong electrolyte in solution} into ions: soluble ionic compounds and strong acids (HCl, HBr, HI, $\ce{HNO3}$, $\ce{HClO4}$, $\ce{H2SO4}$).
\item Keep as written: solids, liquids, gases, weak acids ($\ce{HF}$, $\ce{CH3COOH}$, $\ce{H2CO3}$, \dots), weak bases ($\ce{NH3}$) and water.
\item Cancel the \emph{spectator ions}, which appear unchanged on both sides. Check that atoms and charge balance.
\end{enumerate}
\textbf{Solubility facts the AP exam expects.} All salts of $\ce{Na+}$, $\ce{K+}$ (and the other group 1 ions), $\ce{NH4+}$ and $\ce{NO3-}$ are soluble. Any other solubility information is given in the question.

\worked{2.1} Aqueous iron(III) chloride and sodium hydroxide: $\ce{FeCl3(aq) + 3NaOH(aq) -> Fe(OH)3(s) + 3NaCl(aq)}$. Net: $\ce{Fe^3+(aq) + 3OH-(aq) -> Fe(OH)3(s)}$; spectators $\ce{Na+}$ and $\ce{Cl-}$.

\worked{2.2} Hydrofluoric acid with potassium hydroxide: HF is weak, so $\ce{HF(aq) + OH-(aq) -> F-(aq) + H2O(l)}$, not $\ce{H+ + OH- -> H2O}$.

\errorsh
\begin{itemize}
\item Splitting a precipitate, a weak acid or a gas into ions.
\item Leaving spectator ions in, or canceling ions that actually change.
\item Forgetting that a net ionic equation must balance charge as well as atoms.
\end{itemize}

\section{Representations of reactions (Topic 4.3)}\label{t:4.3}
Particle diagrams must conserve atoms: count every atom of each kind before and after. To find what remains after a reaction in a diagram, treat the particles like moles: find the limiting reactant from the ratio, use it up, and show the leftover excess reactant alongside the products. Spectator ions in solution stay in the ``after'' picture unchanged.

\worked{3.1} A box holds 6 CO and 2 $\ce{O2}$ molecules reacting by $\ce{2CO + O2 -> 2CO2}$. 2 $\ce{O2}$ use 4 CO, so the box afterward holds 4 $\ce{CO2}$ and 2 CO. Check: C $6 = 4 + 2$; O $6 + 4 = 8 + 2$.

\section{Physical and chemical changes (Topic 4.4)}\label{t:4.4}
\begin{itemize}
\item \emph{Physical change}: the substances keep their identity. Phase changes and dissolving molecular solids overcome intermolecular forces only.
\item \emph{Chemical change}: chemical bonds break and form, producing new substances.
\item Dissolving an ionic solid is the gray area: ionic attractions break and ion--dipole attractions form, but no new substance appears. The AP course accepts either classification if you justify it in terms of the bonds or interactions involved.
\end{itemize}
\worked{4.1} Classify (a) electrolysis of water, (b) condensation of steam, (c) dissolving $\ce{KNO3}$. (a) Chemical: O--H bonds break, $\ce{H2}$ and $\ce{O2}$ form. (b) Physical: hydrogen bonds form between intact molecules. (c) Either, justified: ionic attractions are replaced by ion--dipole attractions; the ions are unchanged.

\section{Stoichiometry (Topic 4.5)}\label{t:4.5}
\subsection{The mole road map}
\[ \text{grams A} \xrightarrow{\;\div M_A\;} \text{mol A} \xrightarrow{\;\times\,\text{coef B}/\text{coef A}\;} \text{mol B} \xrightarrow{\;\times M_B\;} \text{grams B} \]
Other entry and exit points: $n = MV$ for solutions, $n = PV/RT$ for gases, $N = nN_A$ for particles.
\subsection{Limiting reactant, excess, yield}
\begin{enumerate}
\item Convert each reactant to moles.
\item Find how much product each could make (or divide each by its coefficient); the smallest wins. That reactant is \emph{limiting}.
\item Theoretical yield comes from the limiting reactant. Excess left = moles present $-$ moles used.
\item Percent yield $= \dfrac{\text{actual}}{\text{theoretical}} \times 100\%$.
\end{enumerate}
\worked{5.1} $\ce{4Fe + 3O2 -> 2Fe2O3}$, with 11.17 g Fe (0.2000 mol) and 4.80 g $\ce{O2}$ (0.150 mol). Fe could make 0.1000 mol $\ce{Fe2O3}$; $\ce{O2}$ could make 0.100 mol. They tie: both are used up exactly, and the yield is $0.100 \times 159.69 = 16.0$ g.

\worked{5.2} $\ce{2C4H10 + 13O2 -> 8CO2 + 10H2O}$, with 5.81 g butane (0.100 mol) and 16.0 g $\ce{O2}$ (0.500 mol). Butane needs $0.100 \times 13/2 = 0.650$ mol $\ce{O2}$; only 0.500 mol is present, so $\ce{O2}$ limits. $\ce{CO2}$: $0.500 \times 8/13 = 0.308$ mol $= 13.5$ g. Butane used: $0.500 \times 2/13 = 0.0769$ mol; left $0.0231$ mol $= 1.34$ g.

\worked{5.3} Gas stoichiometry: $\ce{Zn + 2HCl -> ZnCl2 + H2}$. 0.654 g Zn (0.0100 mol) gives 0.0100 mol $\ce{H2}$, which at 300. K and 0.985 atm occupies $(0.0100)(0.08206)(300.)/0.985 = 0.250$ L.

\errorsh
\begin{itemize}
\item Using grams in the mole ratio.
\item Choosing the limiting reactant by the smaller \emph{mass}.
\item Calculating theoretical yield from the excess reactant.
\end{itemize}

\section{Introduction to titration (Topic 4.6)}\label{t:4.6}
In a titration a solution of known concentration (the \emph{titrant}, in a buret) is added to a measured volume of the \emph{analyte} until the reaction is complete.
\begin{itemize}
\item \emph{Equivalence point}: moles of titrant added are exactly the stoichiometric amount for the analyte.
\item \emph{End point}: the indicator changes color (or a probe signals). A good indicator makes the end point nearly equal to the equivalence point.
\item At equivalence: $M_tV_t \times (\text{ratio}) = M_aV_a$. For $\ce{H2SO4}$ with NaOH the ratio is 1 : 2; for $\ce{Ba(OH)2}$ with HCl it is 1 : 2.
\end{itemize}
\textbf{Error analysis.} Follow the error into $n = MV$.
\begin{center}
\begin{tabular}{p{0.5\linewidth}p{0.4\linewidth}}\toprule
Error & Effect on calculated analyte \\ \midrule
Buret not rinsed with titrant (titrant diluted) & too high (more volume used) \\
Overshooting the end point & too high \\
Air bubble in buret tip that disappears & too high \\
Stopping before the end point & too low \\
Extra water added to the analyte flask & no effect (moles unchanged) \\
Analyte pipet not rinsed with analyte (diluted) & too low \\ \bottomrule
\end{tabular}
\end{center}
\worked{6.1} 15.00 mL of $\ce{Ca(OH)2}$ needs 22.50 mL of 0.0400 M HCl. Moles HCl $= 9.00\times10^{-4}$; $\ce{Ca(OH)2}$ $= 4.50\times10^{-4}$; concentration $= 4.50\times10^{-4}/0.01500 = 0.0300$ M.

\section{Types of chemical reactions (Topic 4.7)}\label{t:4.7}
\begin{center}
\begin{tabular}{p{0.2\linewidth}p{0.36\linewidth}p{0.34\linewidth}}\toprule
Type & What happens & Example \\ \midrule
Acid--base & proton transfer & $\ce{HNO3 + KOH -> KNO3 + H2O}$ \\
Precipitation & ions combine into an insoluble solid & $\ce{Ag+ + Br- -> AgBr(s)}$ \\
Redox & electron transfer (oxidation numbers change) & $\ce{Cu + 2Ag+ -> Cu^2+ + 2Ag}$ \\ \bottomrule
\end{tabular}
\end{center}
Synthesis, decomposition, single replacement and combustion describe the \emph{form} of a reaction; ask whether electrons are transferred to decide if one is also redox. Any reaction with a free element as reactant or product is redox. Combustion of a hydrocarbon is redox.

\worked{7.1} $\ce{2NaHCO3(s) -> Na2CO3(s) + H2O(g) + CO2(g)}$ is a decomposition but not redox: every element keeps its oxidation number (Na $+1$, H $+1$, C $+4$, O $-2$).

\section{Introduction to acid--base reactions (Topic 4.8)}\label{t:4.8}
\begin{itemize}
\item \emph{Brønsted--Lowry acid}: proton ($\ce{H+}$) donor. \emph{Base}: proton acceptor, which needs a lone pair.
\item \emph{Conjugate pair}: two species that differ by one $\ce{H+}$ ($\ce{HF}$/$\ce{F-}$, $\ce{NH4+}$/$\ce{NH3}$, $\ce{H2O}$/$\ce{OH-}$, $\ce{H3O+}$/$\ce{H2O}$).
\item \emph{Amphiprotic} species can do both: $\ce{H2O}$, $\ce{HCO3-}$, $\ce{HSO4-}$, $\ce{H2PO4-}$.
\item Strong acid + strong base: net ionic $\ce{H+ + OH- -> H2O}$. Weak acid or weak base: keep the weak species intact.
\end{itemize}
\worked{8.1} In $\ce{HCN + H2O <=> CN- + H3O+}$: HCN is the acid (conjugate base $\ce{CN-}$); $\ce{H2O}$ is the base (conjugate acid $\ce{H3O+}$).

\worked{8.2} Net ionic for nitrous acid with aqueous ammonia (both weak): $\ce{HNO2(aq) + NH3(aq) -> NO2-(aq) + NH4+(aq)}$.

\section{Oxidation--reduction (redox) reactions (Topic 4.9)}\label{t:4.9}
\subsection{Oxidation numbers}
Free elements: 0. Monatomic ions: their charge. In compounds: F $-1$; O $-2$ (except peroxides, $-1$); H $+1$ with nonmetals; group 1 $+1$, group 2 $+2$. The sum equals the charge of the species.
\subsection{Oxidation and reduction}
\begin{itemize}
\item \emph{Oxidation}: loss of electrons, oxidation number rises. \emph{Reduction}: gain of electrons, oxidation number falls (``OIL RIG'').
\item The species oxidized is the \emph{reducing agent}; the species reduced is the \emph{oxidizing agent}.
\item Electrons lost = electrons gained in a balanced equation.
\end{itemize}
\subsection{Balancing in acid by half-reactions}
\begin{enumerate}
\item Split into oxidation and reduction half-reactions; balance atoms other than O and H.
\item Balance O with $\ce{H2O}$, then H with $\ce{H+}$, then charge with $\ce{e-}$.
\item Multiply so the electrons cancel, add, and simplify.
\end{enumerate}
\worked{9.1} $\ce{Cr2O7^2- + I- -> Cr^3+ + I2}$ in acid. Reduction: $\ce{Cr2O7^2- + 14H+ + 6e- -> 2Cr^3+ + 7H2O}$. Oxidation: $\ce{2I- -> I2 + 2e-}$, times 3. Sum: $\ce{Cr2O7^2- + 14H+ + 6I- -> 2Cr^3+ + 3I2 + 7H2O}$ (6 electrons transferred).

\worked{9.2} Oxidation numbers in $\ce{K2Cr2O7}$: $2(+1) + 2x + 7(-2) = 0$, so Cr is $+6$.

\errorsh
\begin{itemize}
\item Calling the species that is oxidized the oxidizing agent.
\item Balancing atoms but not charge.
\item Giving oxidation numbers per formula instead of per atom.
\end{itemize}

\section{Lab skills and free-response technique}\label{t:lab}
\begin{itemize}
\item \emph{Gravimetric analysis}: precipitate the analyte, filter, wash, dry to constant mass. Incomplete drying makes the precipitate too heavy (result too high); losing solid through the filter makes it too light.
\item \emph{Titration}: rinse the buret with titrant and the pipet with analyte; read the buret to $\pm0.01$ mL at the bottom of the meniscus; add titrant dropwise near the end point.
\item \emph{Reaction stoichiometry labs}: identify the limiting reactant before calculating any yield.
\end{itemize}
In free response, write net ionic equations with phases, show the mole ratio as a step, and for error questions name the measured quantity that changes and the direction it moves the result.

\clearpage
\section*{Summary: Unit 4 on one page}
Use this page the night before. For any row you cannot explain from memory, reread its section and redo its worked example.

\begin{longtable}{@{}p{0.11\linewidth}p{0.40\linewidth}p{0.43\linewidth}@{}}\toprule
\S{} / Topic & Must know & Key equations and tests \\ \midrule
1 / 4.1 & Evidence of reaction; balancing & coefficients only, never subscripts \\[3pt]
2 / 4.2 & Net ionic; spectators; strong vs.\ weak & Na, K, $\ce{NH4+}$, $\ce{NO3-}$ always soluble \\[3pt]
3 / 4.3 & Particle diagrams conserve atoms & limiting reactant in the picture \\[3pt]
4 / 4.4 & Physical vs.\ chemical; dissolving salts & new substance = chemical \\[3pt]
5 / 4.5 & Mole ratio; limiting; excess; \% yield & g $\rightarrow$ mol $\rightarrow$ ratio $\rightarrow$ g \\[3pt]
6 / 4.6 & Equivalence vs.\ end point; error direction & $M_tV_t\,(\text{ratio}) = M_aV_a$ \\[3pt]
7 / 4.7 & Acid--base, precipitation, redox & free element $\Rightarrow$ redox \\[3pt]
8 / 4.8 & Brønsted--Lowry; conjugates; amphiprotic & differ by one $\ce{H+}$ \\[3pt]
9 / 4.9 & Oxidation numbers; agents; half-reactions & e$^-$ lost = e$^-$ gained \\[3pt]
10 / lab & Gravimetric and titration technique & follow each error into $n = MV$ \\ \bottomrule
\end{longtable}

\textbf{Strong acids to memorize.} HCl, HBr, HI, $\ce{HNO3}$, $\ce{HClO4}$, $\ce{H2SO4}$ (first proton). \textbf{Strong bases.} Group 1 hydroxides; $\ce{Ca(OH)2}$, $\ce{Sr(OH)2}$, $\ce{Ba(OH)2}$.
\end{document}
